
\section{Related Work}
\label{sec:related}

\subsection{ML, LLM, and Multi-Agent Phishing Detection}

Traditional ML-based phishing detectors rely on
URL, webpage, and infrastructure features, but
their dependence on predefined features and
training distributions limits generalization
to unseen attacks. Reference-based and
knowledge-based approaches incorporate visual
and contextual evidence to reduce reliance on
historical blacklists
\cite{lin2021phishpedia,liu2023knowledge}.

LLM-based approaches extend detection through
semantic reasoning and tool-assisted investigation.
Wang et al.~\cite{wang2025can} developed an
agentic SMS analyzer that retrieves URL,
domain, screenshot, and brand evidence,
while Li et al.~\cite{li2026phishintentionllm}
employed retrieval-augmented screenshot analysis.
However, single-agent reasoning or pre-acquired
evidence can constrain adaptive investigation
across heterogeneous modalities.

Multi-agent systems distribute analysis among
specialists. Li et al.~\cite{li2025phishdebate}
combined independent modality analysis,
Moderator-guided debate, and Judge-based
adjudication. Chiba et al.~\cite{chiba2026phishlumos}
coordinated specialized agents for campaign-level
mitigation, while Trad et al.~\cite{trad2025clasp}
introduced cost-ordered escalation among URL,
screenshot, and HTML agents. Nevertheless,
agent agreement may reflect correlated evidence,
while verdict-driven escalation and repeated
debate may incur unnecessary costs without
resolving specific evidence gaps.

\subsection{Uncertainty-Aware Reasoning and
Collaborative Verification}

Calibrated uncertainty supports selective
reasoning. Kamath et al.~\cite{kamath2020selective}
studied error estimation under domain shift,
while Chang et al.~\cite{chang2026cascadedebate}
used calibrated confidence for cost-aware
multi-agent cascading. However, confidence-based
routing can degrade under distribution shift
or heterogeneous specialist capabilities
\cite{jitkrittum2023when,ovadia2019can}.

LLM self-reported confidence is also unreliable
for identifying errors~\cite{xiong2024can}.
Repeated sampling improves estimation at
additional inference cost, while logit-free
conformal methods offer alternatives for
black-box models~\cite{su2024api}.
Neither directly establishes evidence-driven
escalation among heterogeneous phishing agents.

Collaborative reasoning introduces agent
deference, whereby agents adopt peer conclusions
without independently examining evidence.
Pitre et al.~\cite{pitre2025consensagent}
addressed debate sycophancy through an
input-side intervention. However, such
interventions do not enforce evidence-supported
revisions. These limitations motivate
record-derived confidence, provenance-aware
collaboration, and output-side validation
rather than reliance on self-reported
confidence or consensus.

\subsection{Zero-Day Evaluation and Research Gaps}

Previous studies evaluated unseen phishing
attacks using forward-in-time testing and
reputation-source exclusion
\cite{lin2021phishpedia,liu2023knowledge}.
For LLM-based detectors, however, these
conditions cannot exclude potential exposure
during model pretraining. Distribution shift
further challenges calibrated routing
\cite{jitkrittum2023when}, highlighting the
need to evaluate both detection and adaptive
collaboration under unseen-attack conditions. A related study, \emph{Confidence-Gated
Multi-agent Phishing Detection}
\cite{li2026phishcascade}, addresses
confidence-guided multi-agent detection.
Its detailed methodology and results were
unavailable for this review, preventing
a mechanism-level comparison.
To the best of our knowledge, MA-ZeroPhish
introduces a novel agentic multi-agent framework
for zero-day phishing detection that enables
adaptive evidence acquisition, provenance-aware
collaborative reasoning, and uncertainty-driven
agent coordination with verifiable
inter-agent revisions.


